\chapter{The universal dictionary and the misfit certificate}\label{ch:dictionary}
\chapterbrief
{a filtered dictionary $\mathcal D_0\subset\mathcal D_1\subset\cdots$ of finite sets of operators (the \emph{atoms}), enumerated by an explicit algorithm; a datum $u$ \emph{fits} a level if some non-zero combination of its atoms annihilates $u$.}
{the project was asked: can a dictionary be built that ``takes everything into account''? The honest answer has three parts: in the sense of approximation, yes; as exact identification of every operator, no; and in between, a misfit can always be certified by finite data, while a fit never can.}
{supplies Steps~2 and~4 of the extended algorithm (Chapter~\ref{ch:extalgo}). The proofs are Dim.~A48--A51.}

\section{What ``universal'' can mean}

\begin{theorem}{Every local linear operator is a limit of dictionary elements \led{DICT-U1} \proved}\label{thm:dictu1}
By Peetre's theorem a linear operator on $C^\infty(\Omega)$ that does not enlarge supports is a differential operator with smooth coefficients, of locally finite order. With a partition of unity and polynomial approximation of the coefficients it is a limit, in the topology of locally bounded order, of elements of the explicit dictionary $\{\chi_j(y)y^\beta\partial^\alpha\}$. It is a pointwise limit of polynomial-coefficient operators iff its order is globally finite.
\dimref{DICT-U1}
\end{theorem}

The same holds for classes beyond differential operators: regular non-linear operators (differential polynomials, by a local version of a Peetre--Slov\'ak argument), symbols with loss $\varepsilon$, completely monotone memory kernels, delays (as an augmented transport system) and systems, each with its own explicit atoms \led{DICT-U2, U3, RIT1, U4}.

\section{What it cannot mean}

\begin{theorem}{No countable dictionary identifies everything \led{DICT-N1} \proved}\label{thm:dictn1}
(i) For any countable dictionary the exactly representable operators form a meagre set (Baire). (ii) No finite data certifies the fit or the identification. (iii) If $u\in\EP$, at every high level $\Ann(u)\cap R_{\le n}\supsetneq(L)\cap R_{\le n}$ for any single law $L$: the kernel is strictly larger than the consequences of one law. (iv) A procedure with finite outputs emits only countably many operators. (v) For a filtered family of differential polynomials the set of $u$ that satisfy some equation of the dictionary is meagre \emph{and} shy (prevalence): the generic $u$ satisfies none.
\dimref{DICT-N1}
\end{theorem}

\begin{intuition}
A compact set in an infinite-dimensional space has empty interior, so a countable family of finite-parameter families covers a ``thin'' set. For the measure statement one moves $u$ along a direction on which the determinant of the dictionary applied to $u+sw$ is a non-zero polynomial in the single real parameter $s$: only finitely many $s$ lie in the bad set.
\end{intuition}

\begin{caution}[title=What this theorem does not say]
``Everything'' holds in the sense of approximation (Theorem~\ref{thm:dictu1}) plus a misfit certificate, not as exact identification. The first form of the statement for coefficients that are $C^\infty$ functions (``no operator is identified'') is false and was replaced by a rank theorem: an affine family is identified iff $\ker\Phi_u=0$ \led{DICT-RANGO}.
\end{caution}

\section{Fit and misfit}

\begin{theorem}{Asymmetry of fit and misfit \led{DICT-A, MF1--MF3} \proved}\label{thm:dicta}
Let $\{D_0,\dots,D_N\}$ be independent continuous operators. The set of $u$ with a \emph{misfit} (the $D_iu$ independent) is open, and each such $u$ has a certificate: a non-vanishing determinant $\det\langle D_iu,\psi_j\rangle$ for finitely many test functions, valid in a neighbourhood. The set of $u$ that \emph{fit} is closed with empty interior, and for differential atoms a fit cannot be certified even from finitely many \emph{exact} measurements.
\dimref{DICT-A}
\end{theorem}

\begin{example}{A certified misfit, an uncertifiable fit}
For $u=e^{-t}\sin x$ the level $\{\partial_t,\partial_x\}$ shows a misfit: the Gram matrix of $(u_t,u_x)$ is non-singular and $\varphi=\langle\cdot,\mathrm{res}\rangle/\|\mathrm{res}\|$ certifies that no first-order law holds. The level $\{\partial_t,\partial_x^2\}$ fits ($u_t=u_{xx}=-u$), and no finite data can certify that.
\end{example}

Misfit certificates exist in three forms: a Hahn--Banach functional (distance to the span), a Gram quotient in a Hilbert space, and a weak form that needs no derivative of $u$ (only integrals against test functions) \led{MF1, MF2, MF3}.

\section{Stopping rules and the limit}

\begin{theorem}{Stopping rule and identification in the limit \led{DICT-GR, I-LIM} \proved}\label{thm:dictgr}
If $n_0$ is the least level with an affine fit and the kernel at every higher level equals the consequences of that law, then the test stops exactly at $n_0$. With a compact class of coefficients and uniqueness, small residual means close coefficients: $\|L_n-L\|\le\omega(C_u\varepsilon_n)$, and with data of error $\delta$ one gets $\omega(C_u\varepsilon_n+2C_{R'}\delta)$, a H\"older rate in the one-coefficient case.
\dimref{DICT-GR}
\end{theorem}

\section{Other results of this phase}

\begin{itemize}
\item \textbf{Exact gain of an atom} \led{DICT-R3}; \textbf{change of frame} by the Volterra encoder, with exact intertwining and constants $\|W\|,\|W^{-1}\|\le1+a^2Q\,e^{a\sqrt Q}$ \led{DICT-R6}.
\item \textbf{Exponential families}: exact polynomialisation, kinks and the $2$-kink \led{DICT-NS}.
\item \textbf{Minimax rate} $\Theta(\varepsilon^{1/3})$ for the potential reconstruction under $M>\delta_0$ \led{DM-$\delta$-MINIMAX}.
\end{itemize}

\begin{takeaway}
A dictionary can be universal for approximation, never for exact identification. Its negative side is both topological and measure-theoretic; its positive side is a one-way certificate: misfit is witnessed by finitely many integrals, fit by none.
\end{takeaway}
